\documentclass{article}
\usepackage{graphicx}
\usepackage{amsmath}
\usepackage{amssymb}
\usepackage{amsfonts}
\usepackage{hyperref}
\usepackage{exercises}
\usepackage{indentfirst}
\usepackage{circledsteps}
\usepackage{amsthm}
\usepackage{csquotes}
\usepackage{yhmath}
\usepackage{comment}
\usepackage{placeins}
\usepackage{geometry}
\usepackage{tikz}
\newcommand*{\oldneg}{\mathord{\sim}}
\newcommand{\Z}{\mathbb{Z}}
\newcommand{\N}{\mathbb{N}}
\let\oldemptyset\emptyset

\theoremstyle{definition}
\newtheorem{definition}{Definition}

\setlength{\parindent}{.5in}
\geometry{margin=0.5in}
\setlength{\parskip}{12pt}
\title{AC Notes}
\begin{document}
\date{}
\maketitle

The form of the axiom of choice I am using here is: Given any non-empty set $S$, there exists a choice function for $S$. The definition of a choice 
function is that: A function $f$ is a choice function of $S$ if $f$ is from $\mathcal{P}(S)-\emptyset$ into $S$. Now, the axiom of choice is not necessary 
to establish the existence of choice functions for non-empty finite sets. For consider some finite set $S$. For any non-empty subset 
$F$ of $S$, by definition, $\exists x(x\in F)$ holds. Also, $\mathcal{P}(S)$ is finite, say it contains $n$ elements (and these elements 
can be arranged in some order $S_1,...,S_n$). Then the formula:
\[ \textrm{(1) }\exists x_1...\exists x_n(x_1\in S_1\wedge ...\wedge x_n\in S_n) \] holds. By the separation principle,
\[ \textrm{(2) }(\forall A_1)...(\forall A_m)(\exists B)(\forall x)[x\in B\equiv\phi(A_1,...,A_m,x)] \] holds for every $m$.  
Now, (1) and (2) together imply that 
\[ \textrm{(3) }\exists x_1...\exists x_n\exists B\forall x[x\in B\equiv \phi(x_1,...,x_n,S_1,...,S_n,x)] \]
where
\[ \phi(x_1,...,x_n,S_1,...,S_n,x) \] is
\[ x=\langle S_1,x_1\rangle\vee ...\vee x=\langle S_n,x_n\rangle \]
Clearly then (3) proves the existence of a choice function for $S$.

The reasoning used in the above cannot be extended to infinite sets, simply because first order logic does not allow formulas to include 
infinitely many variables. Would I be correct in saying that, if we were working in infinitary logic instead of first order logic, the 
existence of choice functions for infinite sets would also follow (and thus the axiom of choice would follow)?

It should be noted that the existence of a choice function for a set $S$ does not mean that we have contstructed a specific choice 
function. Thus the above argument is non-constructive. What this means is that, in saying: every finite set $S$ must have a choice function $f$,
we know nothing more about $f$ other than $f(x)\in x$ for every $x\subseteq S$.

Is the following proof that every denumerable set has a choice function $C$ correct? Suppose $S$ is denumerable. Then, by definition, there is a 
1-1 function from $\omega$ onto $S$. Hence the elements of $S$ can be arranged into a sequence $a_1,a_2,...,a_n,...$. Now consider any subset $F$ 
of $S$. For any $x\in S$, $x=a_i$ for some $i$. Then consider the set $N$ of all numbers $i$ such that $a_i=x$ for 
an $x\in S$. Because $N$ is a set of natural numbers, it has a least element $e$. Then let $C(F)=a_e$.

\end{document}
